\chapter*{Recomendaciones}
\phantomsection
\addcontentsline{toc}{chapter}{Recomendaciones}

A partir de los resultados obtenidos en la investigación, se proponen las siguientes recomendaciones para la evolución y perfeccionamiento del sistema de microservicios:

\begin{enumerate}
    \item \textbf{Migrar la base de datos a SQL Server 2019 o PostgreSQL 15+.} La versión actual (SQL Server 2008) carece de soporte extendido desde 2019. Se recomienda planificar una migración controlada para garantizar seguridad y mantenibilidad a largo plazo, manteniendo la compatibilidad con los microservicios existentes (circuitos-ms, rotaciones-ms).

    \item \textbf{Implementar patrón Event-Driven con NATS JetStream.} Actualmente el sistema usa NATS Core para solicitud-respuesta. Se recomienda migrar a NATS JetStream para implementar eventos asíncronos (ej: \texttt{circuito.actualizado}, \texttt{rotacion.generada}) que permitan a los microservicios reaccionar automáticamente a cambios, mejorando la reactividad del dashboard y la trazabilidad operativa.

    \item \textbf{Extender los microservicios para consumir pronósticos de demanda y datos meteorológicos.} Incorporar fuentes adicionales (pronósticos meteorológicos, curvas de demanda previstas) en un nuevo microservicio \texttt{pronosticos-ms} mejorará la planificación operativa al permitir simulaciones predictivas y rotaciones proactivas.

    \item \textbf{Desarrollar un módulo de simulación de escenarios usando Worker Threads o procesamiento paralelo.} Aprovechar la capacidad multihilo de Node.js para ejecutar múltiples escenarios de rotación en paralelo sin bloquear la interfaz, facilitando la evaluación de diferentes estrategias de afectación. Este módulo podría implementarse como un microservicio separado \texttt{simulador-ms}.

    \item \textbf{Implementar un flujo de integración continua con pruebas automatizadas para todos los microservicios.} Configurar un pipeline (GitHub Actions o similar) que ejecute pruebas unitarias con Jest (backend), pruebas de integración con NATS testing y pruebas E2E con Cypress (frontend) ante cada cambio, asegurando la estabilidad ante futuras actualizaciones y la calidad de los despliegues independientes de cada microservicio.

    \item \textbf{Empaquetar la aplicación como PWA (Progressive Web App).} Dado el diseño responsivo con Tailwind CSS, convertir la SPA en una PWA permitirá su instalación en dispositivos móviles y el acceso sin conexión parcial (con service workers cacheando recursos estáticos), facilitando la supervisión desde puntos remotos de la red local donde la conectividad pueda ser intermitente.

    \item \textbf{Incorporar un modelo de optimización multiobjetivo basado en machine learning.} Utilizar los datos históricos acumulados por la plataforma (en \texttt{ap\_apagones} y registros de rotación) para entrenar modelos que automaticen la ponderación de criterios técnicos, económicos y sociales, minimizando el impacto de las rotaciones en servicios sensibles. Este servicio podría desplegarse como un microservicio independiente \texttt{ml-ms} comunicado vía NATS.

    \item \textbf{Adaptar la arquitectura para su despliegue en otros despachos provinciales.} Estandarizar la API Gateway y parametrizar las configuraciones (circuitos, bloques, prioridades) en el microservicio \texttt{circuitos-ms} para facilitar la réplica del sistema en otras provincias, creando una red interconectada de gestión de déficit a nivel nacional mediante una topología de microservicios federados.
\end{enumerate}